\annexsubsection{\texorpdfstring{希尔伯特--施密特映射 \(^{(2)}\)}{希尔伯特--施密特映射 (2)}}

设 E 是实希尔伯特空间，内积记作 \((x|y)\)。存在从 E 到其对偶的等距同构 \(j_{E}\)，由下式刻画

\[
(x | y) = \langle x, j _ {\mathrm{E}} (y) \rangle \quad \text { for } x, y \text { in } \mathrm{E}\tag{3}
\]

（TVS, V, §1，第~7号，定理 3）。

设 \(E_{1}\) 和 \(E_{2}\) 是两个实希尔伯特空间，u 是从 \(E_{1}\) 到 \(E_{2}\) 的连续线性映射。将从 \(u^{*} = j_{E_{1}}^{-1} \circ {}^{t}u \circ j_{E_{2}}\) 到 \(E_{2}\) 的连续线性映射 \(E_{1}\) 称为 u 的伴随。映射 \(u^{*}\) 由关系刻画

\[
(u (x _ {1}) | x _ {2}) = (x _ {1} | u ^ {*} (x _ {2})) \quad \text { for } x _ {1} \in \mathrm{E} _ {1}, x _ {2} \in \mathrm{E} _ {2}.\tag{4}
\]

如果 \(v\) 是从 \(\mathbf{E}_2\) 到希尔伯特空间 \(\mathbf{E}_3\) 的连续线性映射，则 \((v \circ u)^* = u^* \circ v^*\)。

设 \(E_{1}\) 和 \(E_{2}\) 是两个实希尔伯特空间，u 是从 \(E_{1}\) 到 \(E_{2}\) 的线性映射。由下列公式在 \(E_{1}\) 上定义两个正二次型 H 和 Q

\[
\mathrm{H} (x) = \| x \| ^ {2}, \quad \mathrm{Q} (x) = \| u (x) \| ^ {2} \qquad (x \in \mathrm{E} _ {1}).
\]

命题 3. --- 假设 u 连续。令 \((e_{i})_{i\in I}\) 是 \(E_{1}\) 的正交规范基，\((f_{j})_{j\in J}\) 是 \(E_{2}\) 的正交规范基。则

\[
\operatorname{Tr} (\mathrm{Q} | \mathrm{H}) = \sum_ {i \in \mathrm{I}} \| u (e _ {i}) \| ^ {2} = \sum_ {j \in \mathrm{J}} \| u ^ {*} (f _ {j}) \| ^ {2} = \sum_ {i \in \mathrm{I}} \sum_ {j \in \mathrm{J}} (u (e _ {i}) | f _ {j}) ^ {2}.
\]

对于每个 \(x \in \mathbf{E}_1\)，有 \(\|x\|^2 = \sum_{i \in \mathbf{I}} (x|e_i)^2\)；类似地，对于每个 \(\|y\|^2 = \sum_{j \in \mathbf{J}} (y|f_j)^2\)，有 \(y \in \mathbf{E}_2\)。因此

\[
\begin{array}{r l} \sum_ {i \in \mathrm{I}} \| u (e _ {i}) \| ^ {2} & = \sum_ {i \in \mathrm{I}} \sum_ {j \in \mathrm{J}} (u (e _ {i}) | f _ {j}) ^ {2} \\ & = \sum_ {j \in \mathrm{J}} \sum_ {i \in \mathrm{I}} (e _ {i} | u ^ {*} (f _ {j})) ^ {2} \\ & = \sum_ {j \in \mathrm{J}} \| u ^ {*} (f _ {j}) \| ^ {2}. \end{array}
\]

特别地，数 \(\sum_{i\in I}\| u(e_i)\|^2\) 与 \((e_i)_{i\in I}\) 的正交规范基 \(\mathbf{E}_1\) 的选取无关。

置 \(t = \operatorname{Tr}(\mathbf{Q}|\mathbf{H})\)。按照定义，对于 I 的每个有限子集 \(I'\)，\(I\)

\[
\sum_ {i \in I ^ {\prime}} \| u (e _ {i}) \| ^ {2} = \sum_ {i \in I ^ {\prime}} Q (e _ {i}) \leqslant t,
\]

从而 \(\sum_{i\in I}\| u(e_i)\|^2\leqslant t\)。令 \((e_1',\dots ,e_p')\) 是 \(\mathbf{E}_1\) 中的有限正交规范序列。该序列可以补成 \((e_{\alpha}^{\prime})_{\alpha \in A}\) 的正交规范基 \(\mathbf{E}_1\)。于是

\[
\sum_ {\alpha = 1} ^ {p} \left\| u (e _ {\alpha} ^ {\prime}) \right\| ^ {2} \leqslant \sum_ {\alpha \in \mathrm{A}} \left\| u (e _ {\alpha} ^ {\prime}) \right\| ^ {2} = \sum_ {i \in \mathrm{I}} \left\| u (e _ {i}) \right\| ^ {2}
\]

对所有 \((e_1', \ldots, e_p')\) 取上确界，得到不等式 \(t \leqslant \sum_{i \in I} \| u(e_i) \|^2\)。从而已证明等式 \(t = \sum_{i \in I} \| u(e_i) \|^2\)。

证毕

如果 \(u\) 上的正二次型  是核的，则称 u 是从 \(\mathbf{E}_1\) 到 \(\mathbf{E}_2\) 的希尔伯特--施密特映射。在此情况下，有 \(\mathbf{Q}: x \mapsto \| u(x) \|^2\)\(\mathbf{E}_1\)\(\mathbf{Q} \leqslant \operatorname{Tr}(\mathbf{Q}/\mathbf{H}) \cdot \mathbf{H}\)，因此 u 连续且\(u\)

\[
\| u \| \leqslant \operatorname{Tr} (\mathrm{Q/H}) ^ {1 / 2}.
\]

设 \(u: \mathbf{E}_1 \to \mathbf{E}_2\) 是连续线性映射。由命题 3，u 是希尔伯特--施密特映射，当且仅当存在 \(u\) 的正交规范基 \((e_i)_{i \in \mathrm{I}}\)，使得 \(\mathbf{E}_1\)\(\sum_{i \in \mathrm{I}} \|u(e_i)\|^2 < +\infty\)。在此情况下，\(\mathbf{E}_1\) 的每个正交规范基都具有同一性质。此外，如果 u 是希尔伯特--施密特映射，则由公式 \(u\)（命题 3），其伴随 \(u^*\) 也是希尔伯特--施密特映射。\(\sum_{i \in \mathrm{I}} \|u(e_i)\|^2 = \sum_{j \in \mathrm{J}} \|u^*(f_j)\|^2\)

